\documentclass[11pt,a4paper]{article}
\usepackage[utf8]{inputenc}
\usepackage[T1]{fontenc}
\usepackage{graphicx}
\usepackage{geometry}
\usepackage{hyperref}
\usepackage{booktabs}
\usepackage{xcolor}
\usepackage{titlesec}
\usepackage{caption}

\geometry{margin=1in}
\captionsetup{font=small,labelfont=bf}

\title{A Diachronic Analysis of Gender Neutrality and Demographics Across Task Continuums on wikiHow}
\author{Ebha Baxla \\ School of Liberal Arts (SOLA), Azim Premji University}
\date{April 16, 2026}

\begin{document}

\maketitle

\begin{abstract}
This paper presents a diachronic analysis of gender-based participation and gatekeeping on WikiHow, a major mass-collaborative platform for instructional knowledge. By utilizing a custom-built extraction pipeline and large language models for demographic attribution, we investigate the rigidity of gender roles across diverse "task continuums," ranging from domestic maintenance to occupational professionalization. Our findings reveal persistent asymmetric power dynamics, where established administrative "super-users" systematically police the linguistic modernization introduced by casual contributors. We document the weaponization of neutrality policies and the delegation of gatekeeping to automated agents as a form of "Epistemological Violence," which ultimately shapes the foundational datasets consumed by future artificial intelligence.
\end{abstract}

\section{Introduction}
As digital repositories of human knowledge continue to grow, the mechanisms through which that knowledge is curated, contested, and controlled remain under-researched. WikiHow serves as a unique laboratory for this study due to its specific focus on "how-to" scripts—instructional texts that define the normative boundaries of daily life. This study interrogates whether mass-collaborative platforms act as mirrors of offline inequality or as active foundries that pour and set these disparities into a permanent digital architecture.

\section{Theoretical Framework \& Methodology}
Drawing upon Judith Butler's theory of gender performativity, we argue that every "how-to" step is a performative act. Our methodology utilizes a diachronic crawl of over 11,000 edit summaries from 2005 to 2025. We categorize the wiki into four macro-continuums: \textit{Domestic \& Household}, \textit{Occupational \& Professional}, \textit{Entertainment \& Leisure}, and \textit{Public Policy \& Governance}.

\section{Section A: Demographics \& User Lifecycles}

\begin{figure}[htbp]
\centering
\includegraphics[width=0.7\textwidth]{images/cohort_demographics.png}
\caption{Contributor Demographics by Cohort (2005--2024).}
\label{fig:cohort}
\end{figure}

\begin{figure}[htbp]
\centering
\includegraphics[width=0.6\textwidth]{images/dormancy_decay.png}
\caption{Active Time and Dormancy Decay (Churn Rate) by Gender.}
\label{fig:dormancy}
\end{figure}

\section{Section B: The Task Continuums}

\begin{figure}[htbp]
\centering
\includegraphics[width=0.7\textwidth]{images/longitudinal_continuum.png}
\caption{Longitudinal Continuum Trends: Gender Parity Index across time.}
\label{fig:trends}
\end{figure}

\begin{figure}[htbp]
\centering
\includegraphics[width=0.7\textwidth]{images/gender_flips.png}
\caption{The "Reversal" Articles: Ownership Flips in Battleground topics.}
\label{fig:flips}
\end{figure}

\section{Section C: Gatekeeping, Vandalism \& Toxicity}

\begin{figure}[htbp]
\centering
\includegraphics[width=0.8\textwidth]{images/perp_target_matrix.png}
\caption{Perpetrator-Target Toxicity Matrix crossing gender and article orientation.}
\label{fig:matrix}
\end{figure}

\begin{figure}[htbp]
\centering
\includegraphics[width=0.7\textwidth]{images/gatekeeping_quadrants.png}
\caption{Gatekeeping Quadrant Analysis: Volume vs. Acceptance Rate.}
\label{fig:quadrants}
\end{figure}

\section{Section D: Behavioral \& Ideological Shifts}

\begin{figure}[htbp]
\centering
\includegraphics[width=0.7\textwidth]{images/survival_curve.png}
\caption{Edit Survival Analysis: Probability of survival for gender-neutralizing diffs.}
\label{fig:survival}
\end{figure}

\begin{figure}[htbp]
\centering
\includegraphics[width=0.7\textwidth]{images/chilling_effect.png}
\caption{The "Chilling Effect": Retention Rates Following toxic administrative reverts.}
\label{fig:chilling}
\end{figure}

\section{Conclusion}
Collaborative digital spaces act as active foundries that pour and set disparities into the algorithmic foundation of the future. Authority structures default to prioritizing visual consistency and established hierarchies over equitable knowledge creation. Intervention is required at the level of administrative architecture to prevent these biases from becoming mathematically cement.

\end{document}
